% Início dos Apêndices
\begin{apendicesenv}
\partapendices % Cria uma página separadora (opcional, dependendo do seu template)

% Seu primeiro apêndice (o questionário que já existia)
\chapter{Questionário de Levantamento de Requisitos e Necessidades com Educadores}
\label{apendice:questionario}
Este apêndice apresenta o instrumento de pesquisa (questionário) estruturado e direcionado a professores... 
% (Seu conteúdo do apêndice A fica aqui)

% Seu segundo apêndice (A Portaria do TEiA)
\chapter{Portaria do Projeto TEiA (TED MEC 16162/2025)}
\label{apendice:portaria_teia}

A seguir é apresentada a portaria referente ao Termo de Execução Descentralizada (TED MEC 16162/2025), que fundamenta o Projeto TEiA (Formação Continuada em Inclusão e Suporte para Estudantes com TEA e Neuroatípicos) e institui as diretrizes normativas para o acompanhamento educacional no âmbito do Instituto Federal de Brasília.

% O comando abaixo vai puxar o PDF e anexar no TCC. 
% pages=- (importa todas as páginas da portaria)
% scale=0.85 (reduz o tamanho um pouco para caber o cabeçalho e rodapé do seu TCC)
% pagecommand={} (mantém a numeração de páginas do seu TCC nas folhas do anexo)
\includepdf[pages=-, scale=0.85, pagecommand={}]{portaria_teia.pdf}

\chapter{Questionário de Levantamento de Requisitos e Necessidades com Educadores}
\label{ap:questionario_educadores}

Este apêndice apresenta o instrumento de pesquisa (questionário) estruturado e direcionado a professores e profissionais da educação, com o objetivo de coletar informações sobre as principais dores, desafios na rotina escolar e requisitos funcionais para o suporte ao manejo comportamental de alunos com TDAH e TOD.

\begin{enumerate}
    \item \textbf{Qual é o seu tempo de atuação na área da educação?}
    \begin{itemize}
        \item[ ] Menos de 2 anos
        \item[ ] De 2 a 5 anos
        \item[ ] De 5 a 10 anos
        \item[ ] Mais de 10 anos
    \end{itemize}
    
    \item \textbf{Com que frequência você atua ou já atuou com alunos diagnosticados com TDAH ou TOD em sala de aula?}
    \begin{itemize}
        \item[ ] Frequentemente (rotina diária)
        \item[ ] Ocasionalmente
        \item[ ] Raramente
        \item[ ] Nunca atuei
    \end{itemize}

    \item \textbf{Quais são os maiores desafios encontrados no momento de documentar e registrar crises ou comportamentos atípicos desses alunos? (Selecione mais de uma opção, se aplicável)}
    \begin{itemize}
        \item[ ] Falta de tempo hábil durante o período de aula.
        \item[ ] Formulários institucionais extensos ou complexos.
        \item[ ] Falta de um canal direto e rápido para reportar à coordenação/família.
        \item[ ] Dificuldade em lembrar dos detalhes do ocorrido após o término do turno.
    \end{itemize}

    \item \textbf{No momento em que ocorre uma crise comportamental em sala de aula, você sente que possui fácil acesso a estratégias de manejo imediato baseadas em evidências?}
    \begin{itemize}
        \item[ ] Sim, sempre sei como intervir de forma padronizada.
        \item[ ] Parcialmente, mas sinto falta de suporte proativo ou sugestões rápidas.
        \item[ ] Não, geralmente preciso recorrer à coordenação ou agir por intuição no momento.
    \end{itemize}

    \item \textbf{Como é realizada, atualmente, a comunicação com os familiares para relatar o comportamento e a evolução diária do aluno?}
    \begin{itemize}
        \item[ ] Agenda física escolar escolar.
        \item[ ] Aplicativos de mensagem (como WhatsApp).
        \item[ ] Reuniões periódicas presenciais.
        \item[ ] Não há um fluxo constante de comunicação diária.
    \end{itemize}

    \item \textbf{Em sua visão, o quão útil seria uma aplicação web responsiva (acessível por computador ou celular) que permitisse registrar ocorrências de forma ágil (usando \textit{chips}/seleções rápidas) e que entregasse imediatamente sugestões práticas de manejo pedagógico?}
    \begin{itemize}
        \item[ ] Extremamente útil (otimizaria a rotina).
        \item[ ] Muito útil.
        \item[ ] Indiferente.
        \item[ ] Pouco útil.
    \end{itemize}
\end{enumerate}

\end{apendicesenv}